\scanpage{108}
\deleted{D'autre part, la forme} \inserted{\uline{NB} On verra plus bas comment la forme bil. $\varphi_n$} s'explicite \unsure{en termes de $t$ par $\tau$}.

\textcircled{b}
\[
 \boxed{\unsure{N|n}(1)=1,\qquad n(xy)=n(x)n(y)}
\]
i.e. $n$ est un hom de monoïdes unitaires $A^\times\longrightarrow k^\times$
\inserted{D'où via a), pour $x=1$ $\boxed{t(1)=2}$.}
\deleted{\unsure{Posons}}\deleted{\unsure{Posons}}

\textcircled{\unsure{c|b}}
Posons $\tau(f)=\alpha(t(f))$, et considérons l'application
\[
 \sigma:f\longmapsto\bar f=\tau(f)-f
\]
\deleted{\unsure{On veut que cette application soit}}
On a
\[
 \tau(\bar f)=\tau(f)\underbrace{\tau(1)}_{2}-\tau(f)=\tau(f)
\]
d'où
\[
 (\bar f)\bar{\phantom f}=\tau(\bar f)-\bar f=f
\]
i.e. \deleted{\unsure{on a}}
\[
 \sigma^2=\operatorname{id}_A.
\]
\margin{Notons que par définition $\boxed{\tau(f)\overset{\mathrm{def}}{=}\alpha(t(f))=f+\bar f}$ et on a $\bar1=1$.}
Ceci posé, on veut que l'on ait

\uline{$1^\circ$}) $\sigma$ est un autom. d'algèbre, i.e.
\[
 \boxed{\sigma(xy)=\sigma(x)\sigma(y)}\qquad x,y\in A
\]
ce qui s'explicite ainsi
\[
 2xy=\underset{\substack{\Vert\\\varphi_n(x,\bar y)\ \longrightarrow\ \text{voir d) ci-dessous}}}{\tau(xy)}\inserted{-\tau(x)\tau(y)}+x\tau(y)+y\tau(x)
\]
\margin{\unsure{comb. lin.|cond. lin.} (à coeff. dans $k$) de $1,x,y$}
\uline{NB} Au facteur $2$ près (qui n'est gênant que \deleted{\unsure{en car.}} si $2$ n'est pas inv. dans $k$), on voit donc que la structure d'algèbre de $A$ est connue quand on connaît $A$ en tant que $k$-module, avec l'élément $1_A$ (d'où $\alpha:k\mathrel{-}A$ par $\lambda\mapsto\lambda.1_A$) et les appl.
\[
 t:A\longrightarrow k,\qquad n:A\mathrel{-\!\!-\!\!-}k.
\]
\margin{\uline{NB} Quand $\alpha$ est injectif, les formules \unsure{ci-dessous|ci-dessus} $\alpha n(x)=x\bar x$ \unsure{(de $2^\circ$)} impliquent \textcircled{b} \unsure{ci-dessus} $\begin{gathered}n(1)=1\\n(xy)=n(x)n(y)\end{gathered}$ \unsure{qui}.}

$2^\circ$)
\[
 \boxed{\alpha(\unsure{N|n}(x))=x\bar x}
\]
implique
\[
 \alpha(n(x))=\alpha\bigl(\unsure{\overline{n(x)}|n(\bar x)}\bigr),
\]
\figure{\begin{tikzpicture}\draw (0,0) .. controls (0.35,-0.4) and (0.8,-0.3) .. (0.85,0.2) .. controls (1,1) and (2,1.3) .. (2.4,0.95);\end{tikzpicture}}
\unsure{mais en} \deleted{\unsure{voudra que}} fait
\[
 \deleted{\unsure{3^\circ)}}\boxed{n(\bar x)=n(x)}.
\]
Car
\[
 n(\bar x)=n(\deleted{\unsure{t(x)}}\inserted{\lambda}1-x)
 \underset{\text{cf. \textcircled{a}}}{=}
 n(x)+\deleted{\unsure{t(x)^2}}\lambda^2-\lambda t(x)
 \qquad\text{où }\lambda=t(x)
\]
\inserted{avec $\lambda=t(x)$}
donc on a bien $n(\bar x)=n(x)$.
\mspage{107}
